% Chapter introduction: why the topic arose, who first asked, what it answers.
% STATUS: written 2026-10-21; every claim is cited in the sections it points to.

Chapter~\ref{ch:prime-semiprime} ended with Jordan ideals of a prime ring. This chapter asks what the
Jordan product $a\circ b=ab+ba$ does to the product rule. Herstein defines a Jordan derivation as a
derivation of the Jordan ring of an associative ring, and he credits the idea to Kaplansky
\cite[p.~1104]{Herstein1957}. Every derivation is a Jordan derivation, and Herstein asks about the
converse \cite[p.~1104]{Herstein1957}.

Section~\ref{sec:jordan-type:herstein-theorem} gives his answer for prime rings, with the extension
of Cusack \cite[p.~323]{Cusack1975}. Section~\ref{sec:jordan-type:semiprime} follows the semiprime
case. Brešar proved the semiprime theorem in 1988 \cite[p.~1003]{Bresar1988}. In 2005 he showed
that a Jordan derivation into a bimodule is a derivation on a part of the ring
\cite[p.~411]{Bresar2005}. The implications between the three kinds of map are shown in
Figure~\ref{fig:jordan:implications}.

Three families of maps change the rule. Jordan triple derivations state it for $xyx$
\cite[p.~416]{FosnerIlisevic2008}. Fošner and Ilišević show that on a $2$-torsion free semiprime ring a
generalized Jordan triple derivation is a generalized derivation
\cite[Corollary~4.2, p.~424]{FosnerIlisevic2008}. Generalized Jordan derivations use a second map
\cite[p.~1]{FerreiraGuzzoFerreira2020}. The $(m,n)$-Jordan derivations weight the two terms
\cite[p.~1881]{MajeedOzel2014}. Section~\ref{sec:jordan-type:triple-and-generalized} treats these.

Section~\ref{sec:jordan-type:jordan-on-algebras} looks at particular algebras, among them triangular
algebras \cite[p.~1]{ChengJing2007}. Section~\ref{sec:jordan-type:jordan-homomorphisms} relates the
problem to Jordan homomorphisms, which Brešar and Zelmanov survey \cite[p.~203]{BresarZelmanov2025}.
Section~\ref{sec:jordan-type:star-left} ends the chapter with maps defined by an involution
\cite[p.~416]{FosnerIlisevic2008} and with Jordan left derivations \cite[p.~1]{Vukman1997}.

\hbfigure{04-jordan-type/implications}{Every derivation is a Jordan derivation and every Jordan derivation is a Jordan triple derivation. The dashed converses hold under the conditions shown.}{fig:jordan:implications}
